%%%%%%%%%%%%%%%%%%%%%%%%%%%%%%%%%%%%%%%%%%%%%%%%%%%%%%%%%%%%%%%%%%%%
%        Matematica Libera
%
% Copyright 2022-26 Daniele Zambelli
%
%------------------------------
% Modulo: Insiemi Numerici
%
% Compilation:
% $> lualatex --shell-escape ml01InsiemiNumerici.tex
% $> make4ht ml01InsiemiNumerici.tex -f html5-common_domfilters 
%    -c ml_make4ht -l -u -s -d ./html/mathjax 
%    "mathjax, 3, sec-filename, fn-in, blind" 
% or
% $> ../../../compila.py -l ml01InsiemiNumerici.tex
% $> ../../../compila.py -j ml01InsiemiNumerici.tex
% 
%------------------------------
%
% This work may be distributed and/or modified under the
% conditions of the LaTeX Project Public License, either version 1.3
% of this license or (at your option) any later version.
% The latest version of this license is in
%   http://www.latex-project.org/lppl.txt
% and version 1.3 or later is part of all distributions of LaTeX
% version 2005/12/01 or later.
%
% This work has the LPPL maintenance status `maintained'.
% 
% The Current Maintainer of this work is 
% Daniele Zambelli - daniele.zambelli@gmail.com
%
% This work is part of ``Matematica Libera'' project: 
%
%  https://bitbucket.org/zambu/matematicalibera/src/master/
%
%%%%%%%%%%%%%%%%%%%%%%%%%%%%%%%%%%%%%%%%%%%%%%%%%%%%%%%%%%%%%%%%%%%%

%========================
% Definizione delle directory
%========================
% \newcommand{\basedir}{matematicadolce/}
\newcommand{\depdir}{deposito/}
\newcommand{\magdir}{\depdir magazzino/}
\newcommand{\matdir}{\depdir materiali/}

%========================
% Preambolo
%========================
%------------------------
% Classe del documento  e pacchetti

\documentclass[a4paper,10pt]{\magdir matlibera}
\usepackage{\magdir matliberamat}
\usepackage{\magdir matliberagraf}
% \usepackage{makeidx}
% \usepackage{lipsum}
\makeindex
%------------------------
% Variabili del progetto

\input{\magdir variabili}

%------------------------
% Variabili per questo volume

\newcommand{\editore}{}
\newcommand{\tipo}{Modulo}
\newcommand{\numero}{4}
\newcommand{\titolo}{Equazioni 1}
% \newcommand{\titolodritto}{\tipo \numero: \titolo}
% \newcommand{\pdftitolo}{Insiemi numerici}
\newcommand{\docvers}{\texttt{2.0.0}}
% \newcommand{\edizione}{2023}
% \newcommand{\Edizione}{Edizione}
%\newcommand{\inclusivita}{}%
% {versione accessibile}{versione ciechi}
\newcommand{\anno}{2024}
\newcommand{\mese}{settembre}
\newcommand{\oggi}{\mese\ \anno}
\newcommand{\mcisbn}{}

%------------------------
% Con o senza soluzioni nelle tabelle
% Non tutte le soluzioni sono stampate, ad esempio gli esercizi formati da 
% una tabella da riempire hanno le celle vuote.
% Dove è previsto, è possibile stampare tutte le soluzioni 
% (nel caso precedente le tabelle già completate) commentando la linea con 
% \verb |solallfalse| e decommentando la linea con \verb|solalltrue|.
\newif\ifsolall % va messo nella classe???
\solallfalse    % Non tutte le soluzioni (tabelle vuote)
% \solalltrue     % tutte le soluzioni (tabelle gia compilate)

%========================
% Documento
%========================
\begin{document}
\frontmatter
% Copyright (c) 2023 Daniele Zambelli - daniele.zambelli@gmail.com
% --------------------------------
% Variabili che devono essere istanziate e che vengono usate da 
% frontespizio.
% \newcommand{\serie}{Matematica Libera}
% \newcommand{\tipo}{Modulo} %{Volume}
% \newcommand{\numero}{12}
% \newcommand{\titolo}{Calcolo infinitesimale}
% \newcommand{\inclusivita}{}%{(versione accessibile)}{(versione ciechi)}
% \newcommand{\descr}{Testo a moduli \protect\\ [1em]
%                     per la Scuola Secondaria di $II$ grado}
% \newcommand{\editore}{Pippo}
% \newcommand{\anno}{2023}
% --------------------------------

% Variabili che devono essere istanziate ma che non vengono usate da 
% frontespizio.
% \newcommand{\titolodritto}{\volume \numero: Calcolo infinitesimale}
% \newcommand{\docvers}{\texttt{0.9.9}}
% \newcommand{\mcisbn}{}

\pdfbookmark{frontespizio}{frontespizio}
\thispagestyle{empty}

\newlength{\drop}
\drop=0.15\textheight
\vspace{\drop}
% {\flushleft %\centering
\fontsize{26pt}{0in}%
\selectfont \textbf{\serie}

\vspace{.5em}
\Large \descr
% }

\vspace{1\drop}

\begin{flushright}
\fontsize{18pt}{0pt}%
\selectfont \textbf{\tipo{} \numero}\\[1em]
\fontsize{26pt}{0pt}%
\selectfont \textbf{\titolo}\\[.8em]
% \fontsize{18pt}{0pt}%
\LARGE \textit{\inclusivita}
% \selectfont \textit{\inclusivita}
\end{flushright}

\vspace*{\fill}

\begin{flushright} % \begin{center}
\Large 
\editore{} -- \textbf{\anno} --
\end{flushright} % \end{center}

\normalsize
\newpage

% Copyright (c) 2023 Daniele Zambelli - daniele.zambelli@gmail.com

% --------------------------------
% Variabili che devono essere istanziate e che vengono usate da 
% colophon.
% \newcommand{\serie}{Matematica Libera}
% \newcommand{\titolodritto}{\volume \numero: Calcolo infinitesimale}
% \newcommand{\editore}{Pippo}
% \newcommand{\anno}{2023}
% \newcommand{\coord}{Daniele~Zambelli}
% \newcommand{\autori}{Bruno~Stecca, Daniele~Zambelli}
% \newcommand{\colab}{}
% \newcommand{\texcol}{Claudio~Beccari, Daniele~Zambelli}
% \newcommand{\docvers}{\texttt{0.9.9}}
% \newcommand{\mcisbn}{}
% \pagenumbering{roman}
\thispagestyle{empty}
\pdfbookmark{colophon}{colophon}
\vspace*{\fill}
{\setlength{\parindent}{0em}\small{
% \begin{center}
{\large{\serie\ – \titolodritto}}

Copyright {\textcopyleft} {\anno} \editore
% \end{center}

\affiancati{.70}{.28}{
Questo libro, eccetto dove diversamente specificato, è rilasciato nei 
termini della licenza Creative Commons Attribuzione – Condividi allo stesso 
modo 3.0 Italia (CC BY-SA 3.0) il cui testo integrale è disponibile agli 
indirizzi:
}{
\begin{center} \includegraphics{\magdir img/by-sa.png} \end{center}
}
\url{https://creativecommons.org/licenses/by-sa/3.0/it/}\\
\url{https://creativecommons.org/licenses/by-sa/3.0/it/legalcode}

% Sei libero di:
% riprodurre, distribuire, comunicare al pubblico, esporre in pubblico, 
% rappresentare, eseguire e recitare quest'opera, di modificare quest'opera, 
% alle seguenti condizioni:
% 
% \emph{Attribuzione} --- Devi attribuire la paternità dell'opera nei modi 
% indicati dall'autore o da chi ti ha dato l'opera in licenza e in modo tale 
% da non suggerire che essi avallino te o il modo in cui tu usi l'opera.
% 
% \emph{Condividi allo stesso modo} --- Se alteri o trasformi quest'opera, 
% o se la usi per crearne un'altra, puoi distribuire l'opera risultante solo 
% con una licenza identica o equivalente a questa.

Per maggiori informazioni su questo particolare regime di diritto d'autore 
si legga il materiale informativo pubblicato in:\\
\url{http://www.copyleft-italia.it}.

\vspace{-1em}
\paragraph{Credits:} Il testo deriva da ``Matematica Dolce'':\\ 
\url{https://www.matematicadolce.eu/} \\
che a sua volta è derivato da
``Matematica C3'':\\
\url{https://it.m.wikibooks.org/wiki/Algebra_1} \\
\url{https://it.m.wikibooks.org/wiki/Algebra_2}\\
\url{https://www.larapedia.com/geometria_razionale/
geometria_razionale_nozioni_fondamentali.html}

\vspace{-1em}
\paragraph{Coordinatori del Progetto:} \coord.

\vspace{-1em}
\paragraph{Autori:} \autori\ oltre agli autori di Matematica Dolce.

% \paragraph{Collaboratori:} \colab.

\vspace{-1em}
\paragraph{Progettazione e Implementazione in \LaTeX:} 
          {\texcol}.

\vspace{-1em}
\paragraph{Collaborazione, commenti e suggerimenti:} 
Se vuoi contribuire anche tu al progetto \serie\ o se  vuoi inviare 
commenti o suggerimenti, scrivi a:
\href{mailto:daniele.zambelli@gmail.com}{daniele.zambelli@gmail.com}.

\vspace{2ex}
Versione del documento: {\docvers} del {\oggi}.

% ISBN \mcisbn

\vspace{2ex}
{\scshape{Dati tecnici per l'adozione del libro a scuola}}

Titolo: \serie, \tipo \numero: \titolo\ - \anno.

Codice: \mcisbn 

% Editore: % \href{http://www.matematicamente.it}{\editore}. 

% Anno di edizione: \anno.

% Formato: ebook (\scshape{pdf}).
\cleardoublepage
% \pagebreak
}}

% \newpage
\setcounter{tocdepth}{2}     % Profondità indice
% Rimuove la voce 'indice' dall'indice
\begin{KeepFromToc}
  \tableofcontents
\end{KeepFromToc}

% Decommentare per la prefazione
\chapter{Prefazione}
% \addcontentsline{toc}{chapter}{Prefazione}
% \markboth{Prefazione}{Prefazione}

I contenuti che presentiamo traggono spunto dai materiali pubblicati nel 
corso Matematica Dolce, volumi terzo e quinto. Quei materiali sono stati 
ampiamente revisionati nel contenuto, nella forma con cui vengono 
presentati, nella struttura e nel formato, anche in seguito a un grande 
lavoro di ristrutturazione e semplificazione degli strumenti \LaTeX per 
la produzione del testo.

Più in dettaglio:

\paragraph {Forma}

\begin{itemize} [nosep]
\item 
Avendo usato strumenti \LaTeX diversi la grafica e la pagina 
appaiono diversi.
\item 
Nei margini incominciano ad apparire rimandi a video lezioni.
\item 
I capitoli iniziano con una breve sintesi del contenuto.
\item 
Alla fine di ogni testo è presente un indice analitico.
\end{itemize}

\paragraph {Struttura}~

Il materiale verrà presentato sia 
suddiviso nei 5 anni della scuola di secondo grado sia 
in moduli per argomento.

\paragraph {Formato}~

Oltre al formato pdf orientato alla stampa, il materiale verrà anche 
trasformato nel formato html  con le formule in mathml e mathjax per 
permetterne una più semplice fruizione da parte di chi ha difficoltà a 
usare carta o pdf.

\paragraph {Caratteristiche}~

I cambiamenti apportati giustificano un cambio di nome, 
grati alla ``vecchia'' esperienza di \emph{Matematica Dolce} 
poniamo ora l'accento su un altro aspetto di questo progetto che si 
chiamerà \textbf{Matematica Libera}.

Rimangono inalterate le principali caratteristiche del progetto 
originario:\\[.5em]
{\Large \centering
libertà \quad collaborazione \quad evoluzione \quad polimorfismo \quad
accessibilità.
% \begin{enumerate} [nosep]
% \item libertà,
% \item collaborazione,
% \item evoluzione,
% \item polimorfismo,
% \item accessibilità.
% \end{enumerate}
}

\paragraph {Dove trovarlo}~

L'intero progetto e le ultime versioni dei testi prodotti saranno 
scaricabili da 
\href{https://bitbucket.org/zambu/matematicalibera/downloads/}
{bitbucket.org/zambu/matematicalibera/downloads/}.

\paragraph {Contatti}~

Se questo progetto ti interessa faccelo sapere: 
\href{mailto:daniele.zambelli@gmail.com}{daniele.zambelli@gmail.com}

\vspace{2em}
Un ringraziamento a tutti quelli che useranno e diffonderanno questo 
materiale e
\dots

\dots buon divertimento con la matematica!
\begin{flushright}
Bruno Stecca, Daniele Zambelli
\end{flushright}
% \setcounter{tocdepth}{2}
% \cleardoublepage

%\newpage
%\input{ml01InsiemiNumerici_pref}
%\cleardoublepage
% \newpage
% \newpage

% --------------------------------
% Capitoli
% --------------------------------
\mainmatter
% ..................................................
% Modulo 4: Equazioni 1
\capitolo{\matdir 04_Equazioni/01_Eq_grado_1/eq01/}
         {equazioni}
\capitolo{\matdir 04_Equazioni/01_Eq_grado_1/compl01/}
         {compl_alg_1gra}
%\capitolo{\matdir 04_Equazioni/01_Eq_grado_1/probl01/}
%         {problemi}
\capitolo{\matdir 04_Equazioni/03_Sistemi_lineari/sist01/}
         {sistemi}

% %..................................................
% Azzeramento numerazione capitoli
\renewcommand{\thechapter}{\Alph{chapter}}
\setcounter{chapter}{0}

% generazione indici
\ifdefined\HCode 
.
\else
\cleardoublepage
 \phantomsection
\addcontentsline{toc}{chapter}{\indexname}
 \printindex
\fi
\end{document}
